\documentclass[11pt,letterpaper]{article}
\usepackage{../assets/scientific_report}
\usepackage[brazil]{babel}
\usepackage{amsmath,float,listings,xcolor,geometry,pgfplots}
\pgfplotsset{compat=1.18}
\usetikzlibrary{arrows.meta}
\geometry{headheight=23pt,top=2cm,bottom=2cm,left=2.5cm,right=2.5cm}
\setlength{\parskip}{0.5em}
\setlength{\parindent}{0pt}
\lstset{language=C,basicstyle=\ttfamily\footnotesize,keywordstyle=\color{blue}\bfseries,commentstyle=\color{green!60!black},numbers=left,numberstyle=\tiny\color{gray},numbersep=8pt,frame=single,breaklines=true,showstringspaces=false}

\begin{document}
\pagestyle{plain}
\begin{center}
{\LARGE\bfseries Sobreposição de computação e comunicação}\\[0.4em]
{\normalsize Eugenio Vitor Lopes dos Santos}\\[0.3em]
{\small DCA3703, Programação Paralela, Universidade Federal do Rio Grande do Norte}
\end{center}

\section{Enunciado}
Simular a difusão de calor em uma barra 1D distribuída entre dois ou mais processos MPI, com células extras para trocar as bordas. Comparar três versões: comunicação com \texttt{MPI\_Send} e \texttt{MPI\_Recv}; comunicação não bloqueante com \texttt{MPI\_Isend}, \texttt{MPI\_Irecv} e \texttt{MPI\_Wait}; e uma versão que usa \texttt{MPI\_Test} enquanto calcula os pontos internos.

\section{Fundamentação}
Cada processo armazena um trecho da barra e duas células fantasma, uma de cada lado. A atualização de uma célula usa seu valor e os valores das duas vizinhas no passo anterior:
\[
u_i^{n+1}=u_i^n+0{,}2\left(u_{i-1}^n-2u_i^n+u_{i+1}^n\right).
\]
As células das extremidades locais dependem dos valores recebidos dos processos vizinhos. As células estritamente internas não dependem dessa troca e podem ser atualizadas enquanto ela está pendente. Nas extremidades globais da barra, as células fantasma permanecem em zero.

\begin{figure}[H]
\centering
\begin{tikzpicture}[x=1cm,y=1cm,>=Stealth,font=\small,
    cell/.style={draw,minimum height=9mm,minimum width=13mm},
    ghost/.style={cell,dashed,fill=gray!12},
    edge/.style={cell,fill=orange!20},
    inner/.style={cell,fill=blue!10}]
\node at (0,1.25) {Vizinho esquerdo};
\node at (5.4,1.25) {Memória local do processo: \texttt{current}};
\node at (10.8,1.25) {Vizinho direito};
\node[cell] (left) at (0,0) {$u_m^n$};
\node[ghost] (gleft) at (2,0) {$0$};
\node[edge] at (3.3,0) {$1$};
\node[inner,minimum width=29mm] at (5.4,0) {$2,\ldots,m-1$};
\node[edge] at (7.5,0) {$m$};
\node[ghost] (gright) at (8.8,0) {$m+1$};
\node[cell] (right) at (10.8,0) {$u_1^n$};
\draw[->,thick] (left) -- (gleft);
\draw[->,thick] (right) -- (gright);
\node[align=center] at (2,-1) {Fantasma\\recebida};
\node[align=center] at (5.4,-1) {Interior: pode calcular\\durante a comunicação};
\node[align=center] at (8.8,-1) {Fantasma\\recebida};
\node[align=center,text width=12cm] at (5.4,-2.15)
    {Bordas locais $1$ e $m$: calcular após concluir os recebimentos.\\
     Enviar \texttt{current[1]} à esquerda e \texttt{current[m]} à direita.};
\end{tikzpicture}
\caption{Índices locais, com $m=\texttt{count}$. As setas mostram os recebimentos de um processo com dois vizinhos; cada mensagem contém um \texttt{double}. Nas extremidades globais, o vizinho ausente é \texttt{MPI\_PROC\_NULL} e a célula fantasma permanece em zero.}
\label{fig:halos}
\end{figure}

Na Figura~\ref{fig:halos}, atualizar $i=1$ exige ler \texttt{current[0]}, enquanto atualizar $i=m$ exige ler \texttt{current[m+1]}. Para $2\leq i\leq m-1$, os três valores necessários já pertencem ao processo. Todas as atualizações leem \texttt{current} e escrevem em \texttt{next}; assim, nenhum valor do passo $n$ é substituído por um valor do passo $n+1$ durante o cálculo.

\subsection{Troca bloqueante}
\texttt{MPI\_Send} e \texttt{MPI\_Recv} trocam as bordas antes de qualquer atualização. O código envia primeiro à direita e recebe da esquerda; em seguida, envia à esquerda e recebe da direita. \texttt{MPI\_PROC\_NULL} representa o vizinho ausente nas extremidades globais.

\subsection{Comunicação ponto a ponto não bloqueante}
As operações não bloqueantes retornam sem exigir que a transferência esteja concluída. Enquanto um recebimento estiver pendente, as células fantasma correspondentes não podem ser usadas no cálculo.

\subsubsection{Escondendo latência com sobreposição de computação e comunicação}
Após iniciar a troca, o processo pode atualizar as células internas, que não dependem das bordas recebidas. Só depois da conclusão dos recebimentos atualiza as duas extremidades locais. A versão com \texttt{MPI\_Wait} espera antes de calcular; a versão com \texttt{MPI\_Test} consulta os recebimentos durante o cálculo interno. Essa consulta também tem custo.

\subsubsection{Criação de \texttt{MPI\_Request}}
\begin{itemize}
\item \texttt{MPI\_Isend}: inicia o envio de uma borda e devolve um \texttt{MPI\_Request} para acompanhar a operação.
\item \texttt{MPI\_Irecv}: inicia o recebimento de uma borda e devolve um \texttt{MPI\_Request} para acompanhar a operação.
\end{itemize}

\subsubsection{Consumo de \texttt{MPI\_Request}}
\begin{itemize}
\item \texttt{MPI\_Wait}: espera a conclusão da operação associada ao \texttt{MPI\_Request}.
\item \texttt{MPI\_Test}: verifica se a operação terminou sem bloquear; se ainda estiver pendente, o pedido continua válido.
\end{itemize}

\subsection{Ordem das operações em cada passo}

\begin{figure}[H]
\centering
\begin{tikzpicture}[x=1cm,y=1cm,>=Stealth,font=\footnotesize,
    block/.style={draw,align=center,minimum height=11mm,text width=26mm},
    comm/.style={block,fill=orange!20},
    calc/.style={block,fill=blue!10}]
\draw[->] (0,0.8) -- (13.8,0.8) node[pos=0.5,above] {Ordem de execução (sem escala de tempo)};
\node[anchor=west,font=\small\bfseries] at (0,0) {Send/Recv};
\node[comm] (b1) at (4,0) {Trocar bordas\\Send/Recv};
\node[calc,text width=39mm] (b2) at (8.2,0) {Calcular todas as células\\$1,\ldots,m$};
\node[block,text width=19mm] (b3) at (12.5,0) {Trocar\\ponteiros};
\draw[->] (b1) -- (b2); \draw[->] (b2) -- (b3);
\node[anchor=west,font=\small\bfseries,align=left] at (0,-1.7) {Isend/Irecv\\+ Wait};
\node[comm,text width=19mm] (w1) at (3.5,-1.7) {Iniciar\\4 operações};
\node[comm,text width=19mm] (w2) at (6.1,-1.7) {Wait nos\\4 pedidos};
\node[calc] (w3) at (9.1,-1.7) {Calcular tudo\\$1,\ldots,m$};
\node[block,text width=19mm] (w4) at (12.5,-1.7) {Trocar\\ponteiros};
\draw[->] (w1) -- (w2); \draw[->] (w2) -- (w3); \draw[->] (w3) -- (w4);
\node[anchor=west,font=\small\bfseries,align=left] at (0,-3.6) {Isend/Irecv\\+ Test};
\node[comm,text width=19mm] (t1) at (3.5,-3.6) {Iniciar\\4 operações};
\node[calc] (t2) at (6.5,-3.6) {Calcular interior\\+ Test periódico};
\node[comm] (t3) at (10,-3.6) {Wait se recepção\\ainda pendente};
\node[calc] (t4) at (10,-5.1) {Calcular bordas\\$1$ e $m$};
\node[comm] (t5) at (6.5,-5.1) {Wait nos\\2 envios};
\node[block,text width=19mm] (t6) at (3.5,-5.1) {Trocar\\ponteiros};
\draw[->] (t1) -- (t2); \draw[->] (t2) -- (t3);
\draw[->] (t3) -- (t4); \draw[->] (t4) -- (t5); \draw[->] (t5) -- (t6);
\end{tikzpicture}
\caption{Sequência implementada pelas três versões. Na terceira, seguir as setas até a linha inferior. O cálculo interno oferece uma oportunidade de sobreposição; o diagrama não representa durações medidas nem garante progresso simultâneo das transferências.}
\label{fig:sequencia}
\end{figure}

Na versão com \texttt{MPI\_Wait}, os quatro pedidos são concluídos antes da primeira atualização. Na versão com \texttt{MPI\_Test}, o laço calcula $i=2,\ldots,m-1$ e consulta os recebimentos ainda pendentes quando \texttt{i \% 1024 == 0}. O teste informa a conclusão por uma flag; se ela continuar falsa ao terminar o interior, \texttt{MPI\_Wait} garante que a célula fantasma esteja pronta antes do cálculo das bordas.

Os pedidos \texttt{requests[0]} e \texttt{requests[1]} acompanham os recebimentos; os pedidos \texttt{requests[2]} e \texttt{requests[3]}, os envios. Enquanto um envio estiver pendente, seu buffer em \texttt{current} não pode ser alterado. Escrever os resultados em \texttt{next} permite calcular sem modificar esse buffer. A espera pelos envios ao final libera sua reutilização nos passos seguintes; só então o programa troca os ponteiros dos dois vetores.

\section{Experimento}
O job 2109344 foi executado na partição \texttt{intel-512} do NPAD, com dois processos em nós distintos, uma CPU por processo. O programa foi compilado com \texttt{mpicc -std=c11 -O3 -Wall -Wextra}. A barra tinha 1\,048\,576 células, distribuídas igualmente; cada processo possuía 524\,288 células locais e duas fantasmas. A metade esquerda começou com temperatura 1, e a metade direita com 0. Foram executados 1\,000 passos com coeficiente 0,2, em cinco execuções do programa.

Em cada execução, as três versões foram medidas nessa ordem, com a mesma condição inicial. O intervalo de \texttt{MPI\_Wtime} abrange os passos da simulação e a barreira final; não inclui alocação, inicialização da barra nem redução da soma final. A soma global foi impressa como verificação simples de consistência entre versões. Para comparar os tempos, usa-se a mediana das cinco execuções.

\section{Resultados}
\begin{table}[H]
\centering
\begin{tabular}{crrr}
\toprule
Execução & \texttt{Send/Recv} (s) & \texttt{Isend/Irecv+Wait} (s) & \texttt{Isend/Irecv+Test} (s) \\
\midrule
1 & 0,506098 & 0,498725 & 0,680938 \\
2 & 0,520909 & 0,505166 & 0,682729 \\
3 & 0,486050 & 0,479143 & 0,677186 \\
4 & 0,517211 & 0,490103 & 0,678906 \\
5 & 0,519315 & 0,500133 & 0,677699 \\
\midrule
Mediana & 0,517211 & 0,498725 & 0,678906 \\
\bottomrule
\end{tabular}
\caption{Tempos do job 2109344. A soma final impressa foi 524\,272,535328 em todas as versões e execuções.}
\label{tab:tempos}
\end{table}

\begin{figure}[H]
\centering
\begin{tikzpicture}
\begin{axis}[
    ybar, width=0.9\textwidth, height=6cm,
    ymin=0, ymax=1.8,
    ylabel={Mediana do tempo (s)},
    xtick={1,2,3},
    xticklabels={\texttt{Send/Recv},\texttt{Isend/Irecv+Wait},\texttt{Isend/Irecv+Test}},
    x tick label style={font=\small},
    nodes near coords,
    grid=major,
    bar width=28pt
]
\addplot+[fill=primaryblue] coordinates {(1,0.517211) (2,0.498725) (3,0.678906)};
\end{axis}
\end{tikzpicture}
\caption{Mediana das cinco execuções para cada versão.}
\label{fig:medianas}
\end{figure}

As medianas das versões bloqueante e não bloqueante seguida de espera diferem em aproximadamente 3,6\%, com menor tempo para a versão com \texttt{MPI\_Isend}/\texttt{MPI\_Irecv} e \texttt{MPI\_Wait}. A versão com \texttt{MPI\_Test} teve mediana aproximadamente 31\% maior que a bloqueante: a sobreposição implementada não trouxe ganho neste caso. Uma hipótese é que o custo dos testes periódicos e sua interação com o cálculo superem o tempo de comunicação ocultado; este experimento não isolou esses custos. Cada borda contém apenas um \texttt{double} (8 bytes), portanto a comunicação entre nós é dominada por mensagens pequenas e latência.

As cinco execuções apresentaram valores próximos, sem uma execução discrepante como no teste anterior. A soma final igual é uma verificação global, não uma comparação célula a célula.

\appendix
\clearpage
\section*{Anexo A: Código completo}
\lstinputlisting{tarefa15.c}

\section*{Anexo B: Script Slurm}
\lstinputlisting[language=bash]{tarefa15.slurm}
\end{document}
